\chapter{Discussion and Limitations}
\label{ch:discussion}

\section{Organization of evidence: answering RQ1}
For RQ1, the experiments support combining a broad residual detector with specialists that use different temporal evidence. A specialist is useful when its additional detections justify the errors it introduces into the complete system.

The historical controls and final interventions support different parts of this reasoning. Graph and temporal controls establish that relational context carried predictive information in the episode benchmark. Expert diagnostics expose the limits of family assignment and routing. The later matched interventions assess changes to the selected hybrid. These comparisons explain successive design choices, although they were conducted under different benchmark conditions.

The final branches reflect these choices. General judges broad residual discrepancies, Drift captures directional persistence, and Noise captures excess variability. General's internal tree mixture combines multiple views of a reading; the outer decision process governs the extra alarms supplied by temporal specialists.

Target-group exclusion also gives the reference a clear information constraint: the target reading cannot directly supply its own prediction. Robust fitting, support indicators and temporal context help the detector interpret the resulting discrepancy. Corrupted supporting sensors can still distort the reference. The classifier therefore receives support and reliability information alongside residual magnitude and temporal features. Reference representation and fitting changed together during development, so those historical gains cannot separate their effects. The frozen-rule diagnostics also keep the reference fixed. Section~\ref{sec:reference_factorial} addresses target-group exclusion and iterative reweighting in a separate factorial design with downstream refitting and common calibration constraints. Its comparisons are conditional on the retained representation and learning recipe; whitening and median aggregation remain fixed.

The L-Town history result illustrates this approach within General. Histories of existing residual features failed the first pilot, whereas differences between actually received pressure readings supplied a useful alternative representation. The accepted extension adds these features across the internal mixture and its router. Its full-system evaluation improved replay and all other family means. In this setting, changing what the history represented helped more than extending the history of existing residual features.

The factorial experiment does not show a consistent performance benefit from target-group exclusion. Under its primary calibration policy, iterative reweighting improves pooled F1 in both networks, with a much larger replay contribution in Modena. Exclusion has no consistent pooled-F1 benefit, and its L-Town replay effect changes sign depending on whether iterative reweighting is enabled. The information constraint therefore remains a property of the selected design, while its pooled-F1 benefit is unsupported by this comparison. These results concern refitted detectors under fixed representation and gate settings; they do not identify an independently optimized replacement for the selected system.

\section{Integration and its costs: answering RQ2}
Modena illustrates the precision--sensitivity trade-off in RQ2. The matched integrated candidate raises pooled F1 from 0.6863 to 0.8672 and reduces false positives by 84.2\%, while true positives fall by 4.1\%. Drift and noise F1 decline by 0.0071 and 0.0051. The reduction in false alarms outweighs the lost detections in pooled F1, despite the small drift and noise losses.

This result explains why specialist quality must be assessed at the final alarm level. An accurate specialist can still contribute redundant true positives and many additional false positives. Conversely, suppressing an alarm can improve precision while losing a true detection outside General's coverage. Preserving General protects its existing predictions, but it does not recover every true positive removed from a specialist. Modena shows both the precision benefit and the lost specialist detections.

The threshold-only Modena alternative reaches pooled F1 of 0.6998, well below the integrated candidate. The larger benefit therefore belongs to the package of verification, early-warning context and calibration. A matched factorial ablation would be required to allocate separate effects to those components.

L-Town demonstrates a different balance. Its final history intervention increases pooled F1 from 0.8862 to 0.8950 while improving all family means and reducing mean clean FPR. Small individual drift and targeted losses remain, but the largest other-family loss is approximately 0.0009 F1. The history change brings a smaller improvement than Modena's integration change, but also preserves the required performance of the existing system.

Family macro F1, pooled F1 and clean FPR describe different aspects of these results. The family average describes balanced coverage; the other metrics expose background errors. Modena demonstrates that excellent family scores alone can coexist with weak pooled precision, while the history comparison shows that a weak-family improvement can also benefit the wider detector.

Decision delay is the other explicit cost. Drift and Noise can use evidence through three hours after the endpoint being judged. The historical Modena comparison in Table~\ref{tab:historicaldelay} quantifies the trade-off: drift and noise improve under the delayed configuration, while clean false alarms increase and replay declines slightly. This evidence motivates delayed specialist acceptance, without identifying three hours as an optimal setting for the final architecture. General retains an immediate path, and early warning supplies separate network-time context, but neither removes the delay from the reported specialist task. Whether that allowance is acceptable depends on the response time required by an application. Operational response requires a further evaluation of alarm timing and consequences.

\section{Consistency across networks: answering RQ3}
For RQ3, the repeated fits support stable aggregate performance with distinct family profiles. Modena reaches pooled F1 of 0.8669 and L-Town 0.8950. The original and added five-seed groups produce similar selected-system means, reducing concern that the reported profiles depend on a favorable stochastic fit.

Pooled source-mean SDs are about 0.0118 and 0.0109, compared with model-mean SDs of 0.0034 and 0.0013. The larger variation between sources suggests that further replication should prioritize additional scenario collections. More fits on the existing sources would mainly refine the estimate of training randomness.

The family results also show why both networks are necessary. Modena retains strong replay performance, but drift and noise fall below 0.80 in some source/model cells despite means above that level. L-Town retains strong performance on those specialized families, while replay averages 0.7312 and remains below 0.80 in every cell. The history extension improves replay in all 60 paired comparisons, yet most scores remain below the later 0.75 objective. The replay gain is repeatable, but it remains too small to meet the target in most cells.

L-Town's larger network tests the approach beyond the original development setting, but it is not uniformly more difficult: several family scores exceed their Modena counterparts. Topology, hydraulic behavior, generated observations and fitted representations all differ. For RQ3, the two implementations produce consistent aggregate results, but their weaker families differ and performance still depends on the evaluation source.

\section{What the complexity comparison establishes}
At the original operating points, the hybrid retains all five Modena family means above 0.80 and exceeds the three-hour control calibrated for pooled F1 in pooled performance. Recalibrating that control for family balance recovers drift/noise performance while sacrificing replay and pooled performance. These original comparisons favor the hybrid when temporal-family coverage is the priority.

In L-Town, the same pooled-F1-calibrated control has slightly higher pooled F1 and fewer clean false alarms, while the hybrid has higher replay, drift and noise F1 in all 60 paired cells. Each single classifier uses one decision clock: immediate for the zero-hour control and delayed for the three-hour control. The hybrid preserves an immediate General path alongside delayed specialists. These performance and information-time differences do not quantify runtime or operator workload.

Section~\ref{sec:operating_sensitivity-supp20261008} shows why these original comparisons cannot establish superiority across operating points. Pooled-F1 calibration can reverse both L-Town orderings, giving the hybrid higher pooled F1 but lower replay F1. Under worst-family-F1 calibration, its replay advantage returns at less restrictive caps, with more evaluation false alarms. The component diagnostics also expose a replay cost within the hybrid: retaining General alone or General plus Drift improves L-Town replay while reducing pooled and noise F1. Specialist alarms and verification therefore make measurable, conditional contributions; the simpler classifier remains competitive when their temporal-family gains do not justify the additional errors or structure.

The control shares the retained reference and feature representation and uses a fixed temporal-classifier recipe. The comparison is limited to that recipe. A wider search over joint classifiers could change the comparison. The present experiment reports its two prespecified calibration objectives in full and preserves the original hybrid selection.

External methods such as GDN, TranAD and Graph-MoE remain outside the completed comparison. The thesis prioritizes matched internal controls because its questions concern representation, specialization and alarm integration within the selected framework. Fair external benchmarking would require adaptation and retraining with common endpoint labels, observation masks, training information, calibration conditions and decision-time allowances. Published scores on other datasets cannot supply that comparison. Superiority over those external detectors remains untested.

\section{Limits of the evidence and priorities for further work}
The main external-validity limitation is the observation and attack model. The evaluation uses simulated telemetry, independent missingness and specified corruption processes. These choices make labels and information access inspectable, but they cover only part of the variability that an operational system may encounter. Correlated outages, changing sensor availability, altered demand regimes and attacks adapted to the detector would require separately defined evaluations. Such studies should fix the evaluation protocol before testing and separate the effects of new data conditions from changes to the model.

A second limitation concerns the decision unit. The detector classifies available pressure readings. Clean-period rates of approximately 58.4 and 55.0 false alarms per 100,000 endpoints are informative about background precision, but operator workload depends on how repeated alarms are grouped across sensors and time. A deployment evaluation would therefore need incident aggregation, event-level recall and time to first actionable alert. The selected detectors use pressure inputs. Chapter~\ref{ch:benchmark} explains the indirect role of flow corruption in the stored family codes and clean-period subset.

The replication covers ten fitted models and six evaluation sources per network. These six sources were first used in each network's original five-model confirmation, then reused for the five-seed extension, the ten-model summary and the single-classifier comparison. Operating points were selected on calibration data, but prior knowledge of evaluation outcomes can still influence later study choices. Possible selection bias at the study-design level means that these reused comparisons are not equivalent to an untouched final test. Endpoints within episodes and fits evaluated on the same source remain dependent. Source-level summaries preserve this structure, while larger claims about generalization would benefit from additional independent collections. The supplementary sensitivity protocol was fixed after the original evaluation outcomes were known. Although its thresholds were selected on calibration data before reanalysis, the reused evaluation sources provide retrospective comparisons rather than fresh confirmation. The original selections and research-goal assessment remain in force. The calibration amendment on L-Town must also remain part of the record: the original absolute family-floor constraints were infeasible, and final acceptance follows the declared amended and protected procedures.

Further architectural work should follow the observed errors. In Modena, the priority is recovering drift/noise detections lost through specialist suppression without restoring the clean-period false-positive burden. In L-Town, it is identifying received-observation relationships that separate plausible replay from normal repetition. Any candidate should protect all families and clean periods before a fresh confirmation.

A deployment study would also require computational cost, memory, end-to-end delay and changing-support measurements. Prediction runtime is separate from the three-hour information allowance; the two-network accuracy comparison supplies no runtime scaling estimate. These measurements are needed to assess whether the detector is practical to deploy.
